% Source: https://artofproblemsolving.com/community/c3h1360258_help_please
% Stable identity: aops-topic-1360258-q01
The first line is $y=-\dfrac{1}{2}x-\dfrac{3}{2}$, with slope $-\dfrac{1}{2}$. The second is $y=-\dfrac{a}{3}x-\dfrac{2}{3}$, with slope $-\dfrac{a}{3}$. Non-vertical lines are perpendicular exactly when the product of their slopes is $-1$:
\[
\left(-\dfrac{1}{2}\right)\left(-\dfrac{a}{3}\right)=-1\iff\dfrac{a}{6}=-1\iff a=\boxed{-6}.
\]

\begin{asy}
usepackage("amsmath");
size(6.5cm);
defaultpen(fontsize(10pt));
path ra(pair P,pair u,pair v,real s=0.2){return P+s*unit(u)--P+s*unit(u)+s*unit(v)--P+s*unit(v);}
import graph;
draw((-6.5,0)--(3.6,0),gray,Arrow(4)); draw((0,-3.8)--(0,3.4),gray,Arrow(4));
label("$x$",(3.6,0),S,gray); label("$y$",(0,3.4),W,gray);
pair Pt=(-1/3,-4/3);
draw((-6.5,(6.5-3)/2)--(3,(-3-3)/2),linewidth(0.9));
draw(((-3.5+2/3)/2,-3.5)--((3.2+2/3)/2,3.2),blue+linewidth(0.9));
draw(ra(Pt,(1,-0.5),(1,2),0.45));
dot(Pt);
label("$2y+x+3=0$",(-6.5,1.75),N); label("$3y-6x+2=0$",((3.2+2/3)/2,3.2),E,blue);
label("\mbox{}",(min(currentpicture,user=true)+max(currentpicture,user=true))/2,black); // keeps the text after the figure black when LaTeX typesets the labels inline
\end{asy}

With $a=-6$ the second line is $3y-6x+2=0$, that is $y=2x-\dfrac{2}{3}$, of slope $2$; the lines meet at right angles at $\left(-\dfrac{1}{3},-\dfrac{4}{3}\right)$.
